% !TEX root = ../main.tex
%% Corrigé du sujet type Centrale-Supélec — vérifié par verifications/eval_cc.py.

\corrigeSujetEval{C3}{Sujet type Centrale-Supélec}
\enParties

\begin{corrigeEval}[Généralités]{3 points}{}
\begin{questions}
  \item Les événements $(X_k=i)_{1\leqslant i\leqslant n}$ forment un
    système complet ; la formule des probabilités totales donne
    \[P(X_{k+1}=j)=\sum_{i=1}^{n}P(X_k=i)\,P(X_{k+1}=j\mid X_k=i)=\sum_{i=1}^{n}(\mu_k)_i\,p_{ij}=(\mu_kP)_j\]
    (les termes où $P(X_k=i)=0$ étant nuls par convention). Donc
    $\res{\mu_{k+1}=\mu_kP}$, et par récurrence $\res{\mu_k=\mu_0P^{k}}$.
    \pts{1,5 pt}
  \item Les coefficients de $PQ$ sont des sommes de produits de réels
    positifs ; avec $\mathbf{1}={}^{t}(1,\dots,1)$, « $P$ stochastique »
    signifie $P\geqslant0$ et $P\mathbf{1}=\mathbf{1}$, et
    $(PQ)\mathbf{1}=P(Q\mathbf{1})=P\mathbf{1}=\mathbf{1}$ :
    \res{PQ\text{ est stochastique}}. De même $\mu P\geqslant0$ et
    $(\mu P)\mathbf{1}=\mu\mathbf{1}=1$. \pts{1,5 pt}
\end{questions}
\end{corrigeEval}

\begin{corrigeEval}[Une chaîne à deux états]{4 points}{}
\begin{questions}[start=3]
  \item Calculs directs ; par exemple, $R^{2}$ a pour premier coefficient
    $\frac{b^{2}+ab}{(a+b)^{2}}=\frac{b}{a+b}$, et
    \[R+(1-a-b)S=\frac{1}{a+b}\begin{pmatrix}b+(1-a-b)a&a-(1-a-b)a\\b-(1-a-b)b&a+(1-a-b)b\end{pmatrix}
      =\begin{pmatrix}1-a&a\\b&1-b\end{pmatrix}=P.\]
    \pts{1 pt}
  \item $R$ et $S$ commutent, et leurs produits croisés sont nuls : par
    récurrence (ou par la formule du binôme),
    $\res{P^{k}=R+(1-a-b)^{k}S}$. Si $(a,b)\neq(1,1)$, alors
    $|1-a-b|<1$, donc $P^{k}\to R$ et $\mu_k=\mu_0P^{k}\to\mu_0R$ ; or
    chaque ligne de $R$ vaut $\pi$ et $\mu_0$ est de somme $1$, donc
    $\res{\mu_k\to\pi=\bigl(\tfrac{b}{a+b},\tfrac{a}{a+b}\bigr)}$. Passer à la
    limite dans $\mu_{k+1}=\mu_kP$ donne $\pi P=\pi$. Si $a=b=1$, $P$ échange
    les deux états : $\mu_k$ oscille et ne converge que si
    $\mu_0=\bigl(\frac12,\frac12\bigr)$. \pts{2 pts}
  \item État $1$ = soleil, état $2$ = pluie : $a=0{,}3$, $b=0{,}4$, d'où
    $\pi=\bigl(\frac47,\frac37\bigr)$ : à long terme, la proportion de jours
    ensoleillés est $\res{\dfrac47\approx57\,\%}$. \pts{1 pt}
\end{questions}
\end{corrigeEval}

\begin{corrigeEval}[Le théorème de convergence]{6 points}{}
\begin{questions}[start=6]
  \item Une ligne de $Q$ vaut $1=\sum_jq_{ij}\geqslant n\delta$. Comme
    $\sum_ix_i=0$, $\sum_ix_i\delta=0$, d'où la formule. Les
    $q_{ij}-\delta$ étant positifs :
    \[\lVert xQ\rVert_1\leqslant\sum_j\sum_i|x_i|(q_{ij}-\delta)=\sum_i|x_i|\sum_j(q_{ij}-\delta)
      =\res{(1-n\delta)\lVert x\rVert_1}.\]
    \pts{2 pts}
  \item $\mu_{k+1}-\mu_k=(\mu_k-\mu_{k-1})Q$, et $\mu_k-\mu_{k-1}$ est de
    somme nulle : la question 6 et une récurrence donnent
    $\lVert\mu_{k+1}-\mu_k\rVert_1\leqslant(1-n\delta)^{k}\lVert\mu_1-\mu_0\rVert_1$.
    Pour chaque $j$, la série $\sum_k\bigl((\mu_{k+1})_j-(\mu_k)_j\bigr)$ est
    absolument convergente (comparaison à une série géométrique de raison
    $1-n\delta<1$, ou majorée par $0$ si $n\delta=1$) : la suite
    $\bigl((\mu_k)_j\bigr)$ converge. La limite $\pi$ est positive, de somme
    $1$, et le passage à la limite dans $\mu_{k+1}=\mu_kQ$ donne
    $\res{\pi Q=\pi}$. \pts{2 pts}
  \item Si $\pi$ et $\pi'$ sont stationnaires,
    $\pi-\pi'=(\pi-\pi')Q$ est de somme nulle, donc
    $\lVert\pi-\pi'\rVert_1\leqslant(1-n\delta)\lVert\pi-\pi'\rVert_1$, ce qui
    impose $\pi=\pi'$ si $n\delta>0$ : \res{\text{unicité}}. Enfin
    $\mu_k-\pi=(\mu_0-\pi)Q^{k}$ et $\lVert\mu_0-\pi\rVert_1\leqslant2$ :
    $\res{\lVert\mu_k-\pi\rVert_1\leqslant2(1-n\delta)^{k}}$. \pts{2 pts}
\end{questions}
\end{corrigeEval}

\begin{corrigeEval}[Le classement PageRank]{7 points}{}
\begin{questions}[start=9]
  \item Les pages $1,2,3,4$ ont respectivement $2,1,1,2$ liens :
    \[\res{P=\begin{pmatrix}0&\frac12&\frac12&0\\0&0&1&0\\1&0&0&0\\\frac12&0&\frac12&0\end{pmatrix}},\qquad
      \Bigl(\tfrac25,\tfrac15,\tfrac25,0\Bigr)P=\Bigl(\tfrac25,\tfrac15,\tfrac25,0\Bigr).\]
    La partie C ne s'applique pas : $P$ a des coefficients nuls. La page $4$,
    vers laquelle ne pointe aucun lien, n'est jamais revisitée et reçoit un
    score nul ; plus généralement, une page sans lien sortant, ou un groupe
    de pages se citant entre elles, peut « piéger » l'internaute et rendre
    le classement non unique. \pts{1,5 pt}
  \item Avec probabilité $\alpha$, l'internaute suit un lien ; avec
    probabilité $1-\alpha$, il se « téléporte » sur une page tirée au hasard.
    $G$ est stochastique (combinaison convexe de $P$ et de
    $\frac1nJ$) et ses coefficients sont $\geqslant\delta=\frac{1-\alpha}{n}$,
    avec $1-n\delta=\alpha$ : la partie C donne
    \res{\text{un unique }\pi\text{ stationnaire, et }\lVert\mu_k-\pi\rVert_1\leqslant2\alpha^{k}}.
    \pts{1,5 pt}
  \item La quatrième colonne de $P$ est nulle :
    $\pi_4=(\pi G)_4=\alpha\sum_i\pi_ip_{i4}+\frac{1-\alpha}{4}\sum_i\pi_i=\res{\frac{1-\alpha}{4}}$,
    soit $\frac{3}{80}=0{,}0375$ pour $\alpha=0{,}85$. \pts{1 pt}
  \item \begin{enumerate}[label=\alph*)]
      \item Comme $\sum_i\mu_i=1$ : $(\mu G)_j=\alpha(\mu P)_j+\frac{1-\alpha}{n}$,
        c'est l'expression de \texttt{nu[j]}. Au $p$-ième passage,
        \texttt{mu} et \texttt{nu} représentent $\mu_{p-1}$ et $\mu_p$, et la
        question 7 (avec $\lVert\mu_1-\mu_0\rVert_1\leqslant2$) donne
        $\lVert\mu_p-\mu_{p-1}\rVert_1\leqslant2\alpha^{p-1}\to0$ : la boucle
        se termine. Le test est vrai dès que $2\alpha^{p-1}<10^{-6}$, soit
        $p-1>\frac{14{,}51}{0{,}1625}\approx89{,}3$ : \res{\text{au plus }91\text{ passages}}.
      \item Classement : \res{3,\ 1,\ 2,\ 4}. La page $3$ passe devant la
        page $1$ : elle reçoit des liens de trois pages, dont la page $1$,
        bien classée, alors que la page $1$ n'est citée que par $3$ et $4$.
        On retrouve $\pi_4=0{,}0375$.
      \item Le calcul de \texttt{nu} coûte de l'ordre de $n^{2}$ opérations
        (une somme de $n$ termes pour chacune des $n$ coordonnées), ce qui
        est inenvisageable pour des milliards de pages, alors que la matrice
        est presque entièrement nulle. On parcourt plutôt les liens :
    \end{enumerate}
\begin{python}
def pagerank_creux(succ, alpha, eps):
    n = len(succ)
    mu = [1 / n] * n
    while True:
        nu = [(1 - alpha) / n] * n
        for i in range(n):
            for j in succ[i]:
                nu[j] = nu[j] + alpha * mu[i] / len(succ[i])
        if sum(abs(nu[j] - mu[j]) for j in range(n)) < eps:
            return nu
        mu = nu
\end{python}
    Chaque passage coûte un nombre d'opérations proportionnel à $n$ plus le
    nombre de liens ; avec
    \texttt{succ = [[1, 2], [2], [0], [0, 2]]} (pages numérotées à partir de
    $0$), on retrouve les valeurs de la question b). \pts{3 pts}
\end{questions}
\begin{center}\begin{minipage}{\linewidth}\centering
\begin{tikzpicture}[x=0.22cm,y=0.62cm,line cap=round]
  \foreach \yy/\lab in {0/1,-2/10^{-2},-4/10^{-4},-6/10^{-6}}{
    \draw[otRule,line width=0.3pt] (0,\yy) -- (30.5,\yy);
    \node[font=\scriptsize,anchor=east] at (-0.3,\yy) {$\lab$};}
  \draw[otMuted,line width=0.5pt,-{Stealth[length=1.5mm]}] (0,-7.3) -- (31.5,-7.3);
  \draw[otMuted,line width=0.5pt,-{Stealth[length=1.5mm]}] (0,-7.3) -- (0,0.8);
  \foreach \k in {0,10,20,30} \node[font=\scriptsize,anchor=north] at (\k,-7.35) {$\k$};
  \node[font=\scriptsize,anchor=west] at (31.5,-7.3) {$k$};
  \draw[otNavy,line width=1pt] (0,0.301) -- (30,{0.301-30*0.07058});
  \foreach \k/\v in {0/-0.278,1/-0.803,2/-0.874,3/-1.015,4/-1.386,5/-1.688,6/-1.758,7/-1.899,8/-2.271,9/-2.572,
    10/-2.643,11/-2.784,12/-3.155,13/-3.457,14/-3.527,15/-3.668,16/-4.04,17/-4.341,18/-4.412,19/-4.552,
    20/-4.924,21/-5.225,22/-5.296,23/-5.437,24/-5.808,25/-6.11,26/-6.18,27/-6.321,28/-6.693,29/-6.994,30/-7.065}
    \fill[otOrange] (\k,\v) circle (1.5pt);
  \node[font=\scriptsize,text=otNavy,anchor=south west] at (17,-0.5) {borne $2\alpha^{k}$};
  \node[font=\scriptsize,text=otOrange!85!black,anchor=north east] at (13,-3.8) {$\lVert\mu_k-\pi\rVert_1$};
\end{tikzpicture}\hspace{8mm}%
\begin{tikzpicture}[x=1cm,y=8cm,line cap=round]
  \foreach \p/\v/\lab in {1/0.3797/{0{,}380},2/0.1989/{0{,}199},3/0.3839/{0{,}384},4/0.0375/{0{,}038}}{
    \fill[otGreenDark,rounded corners=0.6pt] ({\p-0.3},0) rectangle ({\p+0.3},\v);
    \node[font=\scriptsize,anchor=south] at (\p,\v) {$\lab$};
    \node[font=\scriptsize,anchor=north] at (\p,-0.005) {page \p};}
  \draw[otMuted,line width=0.5pt] (0.4,0) -- (4.6,0);
\end{tikzpicture}\\*[0.5mm]
\legende{À gauche, l'écart $\lVert\mu_k-\pi\rVert_1$ (échelle logarithmique)
décroît bien plus vite que la borne $2\alpha^{k}$ de la partie C. À droite, le
PageRank des quatre pages pour $\alpha=0{,}85$.}
\end{minipage}\end{center}
\end{corrigeEval}
